% SEGMENT OLP-0352-S01
% Part: lambda-calculus
% Chapter: introduction
% Section: composition

\documentclass[../../../include/open-logic-section]{subfiles}

\begin{document}

\olfileid{lam}{int}{com}
\olsection{\usetoken{S}{lambda definable} சார்புகள் சேர்ப்பின் கீழ் மூடியவை}

\begin{lem}
!!{lambda definable} சார்புகள் சேர்ப்பின் கீழ் மூடியவை.
\end{lem}

\begin{proof}
$f$ ஆனது $h$, $g_0,$ \dots,~$g_{k-1}$ ஆகியவற்றிலிருந்து சேர்ப்பால்
வரையறுக்கப்படுகிறது என்று கொள்க. $h$, $g_0$, \dots,~$g_{k-1}$ ஆகியவற்றை
முறையே $H$, $G_0$, \dots,~$G_{k-1}$ ஆல் !!{lambda defined} ஆக உள்ளன என்ற
கருதுகோளில், ~$f$ ஐ !!{lambda define}s என்று அமையும் ஒரு சொல்~$F$ ஐக்
கண்டறிய வேண்டும். ஆனால்~$F$ ஐ
வெறுமனே பின்வருமாறு வரையறுக்கலாம்:
\[
F(x_0, \dots, x_{l-1}) = H(G_0(x_0, \dots, x_{l-1}),
\dots, G_{k-1}(x_0, \dots, x_{l-1})).
\]
வேறு சொற்களில், லாம்டா கலனத்தின் மொழி சேர்ப்பைப் பிரதிநிதித்துவப்படுத்த
மிகவும் பொருத்தமானது.
\end{proof}

\end{document}
